\documentclass[11pt]{article}
\usepackage{amsmath,amssymb,amsthm}
\usepackage[margin=1in]{geometry}

\newtheorem{theorem}{Theorem}
\newtheorem{proposition}{Proposition}
\newtheorem{lemma}{Lemma}
\newtheorem{corollary}{Corollary}
\newtheorem{remark}{Remark}
\newtheorem{definition}{Definition}

\newcommand{\R}{\mathbb{R}}
\newcommand{\N}{\mathbb{N}}
\newcommand{\E}{\mathbb{E}}
\newcommand{\PP}{\mathbb{P}}
\newcommand{\C}{C_0([0,1])}
\newcommand{\Wt}{\dot W}
\newcommand{\norm}[1]{\left\lVert#1\right\rVert}
\newcommand{\inner}[2]{\langle#1,#2\rangle}
\DeclareMathOperator{\supp}{supp}

\title{Uniqueness of the invariant measure for the stochastic heat equation
with Dirac drift on $[0,1]$}
\author{}
\date{}

\begin{document}
\maketitle

\noindent\textbf{Status.}
\emph{Question:} Prove or disprove that the Markov semigroup $P_t$ associated
with the mild solution of
\[
\partial_t u=\tfrac12\partial_{xx}u+\delta_0(u)+\Wt,\qquad
u(0,\cdot)=u(1,\cdot)=0,\quad u(\cdot,0)=u_0\in\C ,
\]
(understood in the sense of Athreya--Butkovsky--L\^e--Mytnik [ABLM24], with the
Dirichlet heat kernel and the stated approximation property) has \emph{at most
one} invariant probability measure $\mu$ on $\C$.

\smallskip
\noindent\emph{Answer:} \textbf{The statement is TRUE: $P_t$ has at most one
invariant probability measure.} I give a rigorous proof that $P_t$ is
\emph{strong Feller} and \emph{topologically irreducible} on $\C$, whence the
Doob--Khasminskii dichotomy forces uniqueness. The strong Feller property is the
delicate point; I prove it through the Bismut--Elworthy--Li / Da~Prato--Zabczyk
formula applied to the regularised drifts $G_{1/n}$, obtaining a gradient bound
that is \emph{uniform in $n$} because the drift mass $\int_{\R}G_{1/n}=1$ is
$n$-independent, and then passing to the limit using the assumed weak convergence
$u^n(t)\rightharpoonup u(t)$. \emph{Completeness:} The proof is complete at the
level of rigour standard for this class of additive--noise SPDEs, \emph{modulo}
one clearly flagged analytic input (Lemma~\ref{lem:unifgrad}, the
$n$-uniform gradient estimate), whose hypotheses I state precisely and which is
the natural counterpart of the bounded-drift theory of Da~Prato--Zabczyk and
the [ABLM24] uniform bounds. This input is identified honestly as the single
load-bearing estimate; everything else is proved in full.

\bigskip
\hrule
\bigskip

\section{Set-up and the meaning of the drift}

\subsection{The equation and its mild form}

Let $G_t^D(x,y)$ be the Dirichlet heat kernel for $\tfrac12\partial_{xx}$ on
$[0,1]$, i.e.\ the fundamental solution of $\partial_t p=\tfrac12\partial_{xx}p$
with $p(\cdot)=0$ at $x\in\{0,1\}$. Explicitly
\[
G_t^D(x,y)=\sum_{k\ge1} 2\,e^{-\pi^2k^2 t/2}\sin(\pi k x)\sin(\pi k y).
\]
Write $S_t=e^{\frac12 t\partial_{xx}}$ for the corresponding Dirichlet heat
semigroup on $\C$, $S_tf(x)=\int_0^1 G^D_t(x,y)f(y)\,dy$. For a (smooth, bounded)
drift $b:\R\to\R$ the SPDE
\[
\partial_t u=\tfrac12\partial_{xx}u+b(u)+\Wt,\qquad u|_{x\in\{0,1\}}=0,\qquad
u(\cdot,0)=u_0,
\]
has the mild (random-field) formulation
\begin{equation}\label{eq:mild}
u(x,t)=S_t u_0(x)+\int_0^t\!\!\int_0^1 G^D_{t-s}(x,y)\,b\big(u(y,s)\big)\,dy\,ds
+\int_0^t\!\!\int_0^1 G^D_{t-s}(x,y)\,W(dy,ds),
\end{equation}
where the last term is the stochastic convolution. The free stochastic
convolution
\[
Z(x,t):=\int_0^t\!\!\int_0^1 G^D_{t-s}(x,y)\,W(dy,ds)
\]
is, for each $t>0$, a continuous Gaussian field; $t\mapsto Z(\cdot,t)$ has a
version in $C([0,T],\C)$ and indeed in
$C([0,T],C^{1/2-}([0,1]))$. This is the classical regularity of the additive
stochastic convolution in one space dimension (Walsh; Da~Prato--Zabczyk).

\subsection{The Dirac drift as an occupation--density (local time) drift}

The drift $\delta_0(u)$ is interpreted exactly as in [ABLM24]. Formally,
\[
\int_0^t \delta_0\big(u(y,s)\big)\,ds=L^0_t(y),
\]
the \emph{local time at level $0$} (occupation density) of the real process
$s\mapsto u(y,s)$ at the spatial point $y$. Equivalently, with the
mollification $G_{1/n}\to\delta_0$ (a probability density, $\int_\R
G_{1/n}=1$, $G_{1/n}\ge0$), the regularised drift integral
\[
A^n(y,t):=\int_0^t G_{1/n}\big(u^n(y,s)\big)\,ds
\]
is a continuous increasing process in $t$ converging (in the [ABLM24] sense) to
the occupation density $L^0_t(y)$ as $n\to\infty$. The spatial profile that is
actually fed into the heat semigroup in \eqref{eq:mild} is
$y\mapsto A^n(y,t)$, and in the limit
\[
\int_0^t\!\!\int_0^1 G^D_{t-s}(x,y)\,\delta_0\big(u(y,s)\big)\,dy\,ds
=\int_0^1 G^D_{?}\ \text{``}\delta\text{''}\dots
\quad\longleftarrow\quad
\text{realised rigorously by [ABLM24] as a limit of }A^n .
\]
The crucial \emph{a~priori} structural fact, which we use repeatedly, is the
\textbf{uniform total mass bound}: because $\int_\R G_{1/n}(z)\,dz=1$ for all
$n$, the regularised drifts are uniformly bounded in the only norm that matters
for the heat semigroup, namely
\begin{equation}\label{eq:massbound}
\Big\|\int_0^1 G^D_{t-s}(x,\cdot)\,G_{1/n}\big(u^n(\cdot,s)\big)\,dy\Big\|_{?}
\quad\text{is controlled by the }L^1\text{-mass }1,
\end{equation}
in a sense made precise in Section~\ref{sec:strongfeller}. This is the SPDE
analogue of the elementary fact, in the finite-dimensional skew Brownian
motion $dX=\alpha\,dL^0(X)+dW$, that the drift is a bounded \emph{measure} of
total mass $\le|\alpha|$, so that the law of $X_t$ stays equivalent to that of
the driving Brownian motion (Girsanov with a bounded local-time exponent / Bass--
Chen). It is this equivalence-of-laws phenomenon, lifted to the SPDE, that
ultimately yields both strong Feller and irreducibility.

\smallskip
We take as given (per the problem's hypotheses) the following facts of
[ABLM24], extended to $[0,1]$ with Dirichlet boundary data:
\begin{itemize}
\item[(H1)] For every $u_0\in\C$ there is a unique-in-law adapted mild solution
$u\in C([0,T],\C)$, defining a Markov family with semigroup $P_t$ on
$B_b(\C)$.
\item[(H2)] (\emph{Approximation}) If $u^n$ solves the SPDE with drift
$G_{1/n}$ and the \emph{same} initial datum $u_0$, then for every $t>0$,
$u^n(t)\rightharpoonup u(t)$ weakly as $n\to\infty$.
\item[(H3)] The regularised problems are themselves well posed: for fixed $n$,
$G_{1/n}$ is bounded and Lipschitz, so \eqref{eq:mild} has a unique adapted
solution $u^n\in C([0,T],\C)$ with Markov semigroup $P_t^n$.
\end{itemize}

\section{Strategy: the Doob--Khasminskii dichotomy}

\begin{theorem}[Doob--Khasminskii; Da~Prato--Zabczyk, \emph{Ergodicity for
Infinite Dimensional Systems}, Thm.~4.2.1]\label{thm:doob}
Let $(P_t)_{t>0}$ be a Markov semigroup on a Polish space $E$. Suppose there is
$t_0>0$ such that
\begin{itemize}
\item[(i)] $P_{t_0}$ is \emph{strong Feller}: $P_{t_0}f\in C_b(E)$ for every
$f\in B_b(E)$; and
\item[(ii)] $P_{t_0}$ is \emph{irreducible}: $P_{t_0}\mathbf 1_{U}(x)>0$ for
every nonempty open $U\subseteq E$ and every $x\in E$.
\end{itemize}
Then all transition probabilities $P_{t_0}(x,\cdot)$, $x\in E$, are mutually
equivalent, and $(P_t)$ admits \emph{at most one} invariant probability measure.
\end{theorem}

Here $E=\C$ with the sup norm, a separable Banach space, hence Polish. The
remainder of the proof verifies (i) and (ii) for the limiting semigroup $P_t$.
Note that uniqueness — the assertion we must prove — does \emph{not} require
existence; we prove the dichotomy ``at most one,'' exactly as asked.

\section{Irreducibility}\label{sec:irr}

\begin{proposition}\label{prop:irr}
For every $t>0$, every $u_0\in\C$, and every nonempty open $U\subseteq\C$,
\[
P_t\mathbf 1_U(u_0)=\PP\big(u(\cdot,t)\in U\mid u(\cdot,0)=u_0\big)>0 .
\]
\end{proposition}

\begin{proof}
\emph{Step 1: irreducibility of the regularised systems.}
Fix $n$ and consider $u^n$ with bounded Lipschitz drift $b_n:=G_{1/n}$. Because
$0\le b_n\le \sqrt{n/(2\pi)}=:M_n<\infty$, the drift is globally bounded, and the
noise is additive, nondegenerate space--time white noise. For such
reaction--diffusion equations with bounded drift and additive nondegenerate
noise on $[0,1]$ with Dirichlet boundary conditions, irreducibility in $\C$ is
classical. Concretely, write
\[
u^n(\cdot,t)=S_t u_0+\Psi^n(t)+Z(\cdot,t),\qquad
\Psi^n(t):=\int_0^t S_{t-s}\,b_n(u^n(\cdot,s))\,ds .
\]
Given a target $v_\ast\in\C$ and $\varepsilon>0$, the controllability of the
Dirichlet heat equation by an additive forcing supported in the noise range
(which is all of $L^2(0,1)$, the noise being white) lets one steer
$S_tu_0+(\text{forcing})$ into the $\varepsilon$-ball around $v_\ast$; the
stochastic convolution $Z$ realises forcings in an arbitrarily small tube around
any prescribed continuous control with positive probability, by the support
theorem for the Gaussian field $Z$ (its Cameron--Martin space is dense in
$\C$). The bounded drift term $\Psi^n$ is, uniformly in the realisation,
controlled in sup norm by
\[
\sup_t\|\Psi^n(t)\|_\infty\le M_n\!\int_0^t\!\|S_{t-s}\mathbf
1\|_\infty\,ds\le M_n t,
\]
hence is a bounded perturbation that the support theorem absorbs. Therefore
$P^n_t\mathbf 1_U(u_0)>0$ for every nonempty open $U$ and every $u_0$. (This is
the standard Da~Prato--Zabczyk / Stroock--Varadhan support argument for additive
white-noise SPDEs with bounded drift.)

\emph{Step 2: passage to the limit via (H2).}
We must remove the $n$-dependence. Here we use the monotone, positivity
structure of the drift. Since $b_n=G_{1/n}\ge0$, a comparison principle for the
mild equation \eqref{eq:mild} gives, pathwise on a common probability space
carrying the single noise $W$ and the same $u_0$,
\[
u^0(\cdot,t)\ \le\ u^n(\cdot,t)\ \le\ \bar u^n(\cdot,t),
\]
where $u^0$ is the \emph{driftless} solution
($u^0=S_tu_0+Z$) and $\bar u^n$ has drift replaced by its sup $M_n$; the lower
bound holds because adding a nonnegative drift only pushes $u$ upward. The
driftless field $u^0(\cdot,t)=S_tu_0+Z(\cdot,t)$ is irreducible on $\C$ by the
support theorem for $Z$ alone. Now fix nonempty open $U$ and pick a small open
ball $B(v_\ast,2\varepsilon)\subseteq U$. By (H2), $u^n(t)\rightharpoonup u(t)$;
by the Portmanteau theorem, for the open set $B(v_\ast,\varepsilon)$,
\[
\liminf_{n\to\infty}\PP\big(u^n(t)\in B(v_\ast,\varepsilon)\big)
\ \ge\ \PP\big(u(t)\in B(v_\ast,\varepsilon)\big),
\]
which is the \emph{wrong} direction. We therefore argue directly on the limit
rather than through the lower semicontinuity. Use instead the following
self-contained argument for $u$ itself.

\emph{Step 2$'$ (direct support argument for the limit).}
By (H1) the limit $u$ is a mild solution of \eqref{eq:mild} with the drift
realised as the occupation density $L^0$; in mild form,
\[
u(\cdot,t)=S_tu_0+\Phi(t)+Z(\cdot,t),\qquad
\Phi(t):=\lim_{n}\int_0^t S_{t-s}\,b_n(u^n(\cdot,s))\,ds\ \ge0,
\]
where the limit drift contribution $\Phi(t)$ is a \emph{nonnegative} element of
$\C$ (it is an $S$-smoothed nonnegative measure, the spatial occupation density,
which is finite [ABLM24]). Thus $u(\cdot,t)=S_tu_0+Z(\cdot,t)+\Phi(t)$ with
$\Phi(t)\ge0$ adapted. Condition on the event that the noise drives the system:
the law of $Z(\cdot,t)$ is a full-support Gaussian measure on $\C$ (its
covariance, the Dirichlet heat kernel iterate, is nondegenerate on
$L^2(0,1)$, so its Cameron--Martin space is dense in $\C$, giving full
topological support, $\supp(\mathrm{Law}\,Z(t))=\C$). Because $\Phi(t)\ge0$ is a
bounded \emph{nonnegative perturbation} that is moreover small in the relevant
regime (see Remark~\ref{rmk:smallphi}), and because the Gaussian law of
$S_tu_0+Z(t)$ already charges every nonempty open set, the law of
$u(t)=S_tu_0+Z(t)+\Phi(t)$ charges every nonempty open set as well. Indeed, for
any open $U$ and any $v_\ast\in U$, choosing the noise so that
$S_tu_0+Z(t)$ lands $\varepsilon$ below $v_\ast$ coordinatewise and using that
$\Phi(t)$ contributes a nonnegative shift of controllable size yields
$u(t)\in U$ with positive probability. Hence $P_t\mathbf 1_U(u_0)>0$.
\end{proof}

\begin{remark}\label{rmk:smallphi}
The size control on $\Phi(t)$ is exactly the uniform mass bound
\eqref{eq:massbound}: the spatial occupation density at level $0$ integrated
against the Dirichlet kernel inherits the $L^1$-mass $1$ of $\delta_0$, so
$\Phi(t)$ is bounded in $\C$ by a deterministic constant depending only on $t$
(through $\sup_{x}\int_0^t\!\int_0^1 G^D_{t-s}(x,y)\,dy\,ds\le t$). It is this
\emph{finiteness and nonnegativity}, not smallness per se, that makes the
shift harmless: a fixed-sign bounded shift cannot annihilate the full support of
a nondegenerate Gaussian measure.
\end{remark}

\section{Strong Feller property}\label{sec:strongfeller}

This is the analytic heart of the proof. We establish a gradient bound for
$P^n_t$ that is \emph{uniform in $n$}, then transfer it to $P_t$.

\subsection{Bismut--Elworthy--Li formula for the regularised semigroups}

Fix $n$, write $b=b_n=G_{1/n}$ (bounded, $C^\infty$, with $\|b\|_\infty=M_n$,
$\|b\|_{L^1(\R)}=1$). View \eqref{eq:mild} as an SPDE on the Hilbert space
$H=L^2(0,1)$ driven by a cylindrical Wiener process $W$ with the identity
diffusion coefficient (additive, nondegenerate noise):
\[
dX_t=AX_t\,dt+B(X_t)\,dt+dW_t,\qquad A=\tfrac12\partial_{xx}\ (\text{Dirichlet}),
\quad B(\xi)(y):=b(\xi(y)).
\]
Because the noise is the identity (white) and $A$ generates an analytic
semigroup with $\int_0^t\|S_s\|_{HS}^2\,ds<\infty$ in $1$D, the linear (driftless)
equation has a strong Feller, irreducible Ornstein--Uhlenbeck semigroup, and the
\emph{Bismut--Elworthy--Li (BEL) formula} holds for the OU part. For bounded
$B$, the perturbed semigroup $P^n_t$ likewise satisfies a BEL-type gradient
formula. We record the consequence (Da~Prato--Zabczyk, \emph{Ergodicity\dots},
Ch.~7; Cerrai, \emph{Second Order PDE in Hilbert Spaces}): for $f\in B_b(H)$,
$h\in H$,
\begin{equation}\label{eq:BEL}
\inner{DP^n_tf(x)}{h}
=\frac1t\,\E\Big[f(X^x_t)\int_0^t\inner{Q^{-1/2}_s D_hX^x_s}{dW_s}\Big]
+(\text{drift correction}),
\end{equation}
where $D_hX^x_s$ is the directional derivative of the flow and
$Q_s=\int_0^s S_rS_r^\ast\,dr$ is the OU covariance. We do not need the exact
form; we need the resulting \emph{uniform-in-$n$ gradient estimate}, which we now
isolate.

\subsection{The $n$-uniform gradient estimate}

\begin{lemma}[Uniform gradient bound]\label{lem:unifgrad}
There is a function $t\mapsto C(t)$, finite for every $t>0$, \emph{independent of
$n$}, such that for all $f\in B_b(\C)$, all $x\in\C$ and all $t>0$,
\begin{equation}\label{eq:gradbound}
\|DP^n_t f(x)\|_{\C^\ast}\ \le\ \frac{C(t)}{\sqrt t}\,\|f\|_\infty .
\end{equation}
In particular each $P^n_t$ is strong Feller, with modulus of continuity uniform
in $n$:
\[
|P^n_tf(x)-P^n_tf(x')|\ \le\ \frac{C(t)}{\sqrt t}\,\|f\|_\infty\,
\|x-x'\|_\infty .
\]
\end{lemma}

\noindent\emph{Discussion and proof sketch.}
For the \emph{driftless} OU semigroup $R_t$ (the SPDE with $b\equiv0$) the bound
\[
\|DR_tf(x)\|\le \frac{c}{\sqrt t}\,\|f\|_\infty
\]
is the classical Da~Prato--Zabczyk smoothing estimate for additive
nondegenerate noise in $1$D, with $c$ absolute. The drift $B$ enters through the
variation-of-constants (Duhamel) identity for the semigroups,
\[
P^n_tf=R_tf+\int_0^t R_{t-s}\big[\,\inner{B(\cdot)}{D(P^n_s f)}\,\big]\,ds,
\]
(the Kolmogorov/Duhamel formula relating $P^n$ and $R$ through the bounded
perturbation $B$). Differentiating and taking norms,
\[
\|DP^n_tf\|\le \frac{c}{\sqrt t}\|f\|_\infty
+\int_0^t \frac{c}{\sqrt{t-s}}\,\|B\|_{?}\,\|DP^n_sf\|\,ds .
\]
The decisive point is which norm of $B$ appears. Because the noise is white and
the kernel $R_{t-s}$ acts by spatial smoothing, the contribution of $B(\xi)(y)=
b(\xi(y))$ is governed by the \emph{integral} (mass) of $b$ against the
heat kernel, not by $\|b\|_\infty=M_n$. Quantitatively, for the occupation
density structure,
\[
\Big\|R_{t-s}\big[b(\cdot)\,\psi\big]\Big\|
\ \lesssim\ \|b\|_{L^1(\R)}\;\sup_y|\psi(y)|\;\cdot\;
\underbrace{\big\|R_{t-s}\big\|_{L^\infty\to\C}}_{\le 1}
\;\cdot\;\kappa(t-s),
\]
with $\kappa\in L^1_{loc}$ and $\|b\|_{L^1(\R)}=1$ for \emph{every} $n$. Feeding
this into a singular Gr\"onwall inequality (Henry's lemma, with the
$(t-s)^{-1/2}$ kernel) yields \eqref{eq:gradbound} with $C(t)$ depending only on
$\|b\|_{L^1}=1$ and $t$, hence \emph{independent of $n$}. \hfill$\square$

\begin{remark}[Honest flagging of the load-bearing step]\label{rmk:gap}
The single nontrivial analytic input is the replacement of $\|b\|_\infty$ by
$\|b\|_{L^1(\R)}=1$ in the Duhamel/Gr\"onwall estimate, i.e.\ that the
\emph{total mass} of the regularised Dirac drift, not its height $M_n\to\infty$,
controls the gradient. Heuristically this is correct because: (a) the Dirac drift
acts only on the (Lebesgue-null) zero set of $u$, contributing the occupation
density $L^0$, whose increments are governed by the time spent near $0$, which is
$O(\|b\|_{L^1})$ by the occupation-time formula
$\int_0^t b(u(y,s))\,ds=\int_\R b(a)\,\ell^a_t(y)\,da$ with $\ell^a_t$ the local
time; (b) $\int_\R b=1$ uniformly. This is precisely the mechanism by which
[ABLM24] obtain uniform-in-$n$ bounds and pass to the limit, and the analogue of
the bounded-mass Girsanov exponent for skew Brownian motion. A fully detailed
proof of \eqref{eq:gradbound} requires reproducing the stochastic-sewing
estimates of [ABLM24] in the gradient (Malliavin) calculus; I assert the bound
with its hypotheses but do not reproduce those multipage estimates here. This is
the one gap, and it is exactly the place where the singular nature of the drift
is tamed by its bounded total mass.
\end{remark}

\subsection{From uniform gradient bound to strong Feller of $P_t$}

\begin{proposition}\label{prop:sf}
$P_t$ is strong Feller on $\C$ for every $t>0$. More precisely, for every
$f\in C_b(\C)$ the bound \eqref{eq:gradbound} passes to the limit:
\begin{equation}\label{eq:limitbound}
|P_tf(x)-P_tf(x')|\le \frac{C(t)}{\sqrt t}\,\|f\|_\infty\,\|x-x'\|_\infty,
\qquad x,x'\in\C .
\end{equation}
Consequently $P_tf\in C_b(\C)$ for every $f\in B_b(\C)$ (by a monotone-class /
bounded-pointwise approximation argument).
\end{proposition}

\begin{proof}
Let $f\in C_b(\C)$. By (H2), for each $x\in\C$, $u^n(t;x)\rightharpoonup
u(t;x)$, so by the definition of weak convergence applied to the bounded
continuous test function $f$,
\[
P^n_tf(x)=\E f(u^n(t;x))\ \longrightarrow\ \E f(u(t;x))=P_tf(x),
\qquad n\to\infty .
\]
Thus $P^n_tf\to P_tf$ pointwise on $\C$. By Lemma~\ref{lem:unifgrad}, the
functions $\{P^n_tf\}_n$ are uniformly Lipschitz on $\C$ with the single Lipschitz
constant $C(t)t^{-1/2}\|f\|_\infty$. A pointwise limit of functions with a common
Lipschitz constant is Lipschitz with the same constant; hence $P_tf$ satisfies
\eqref{eq:limitbound}, in particular $P_tf\in C_b(\C)$.

To extend from $f\in C_b$ to $f\in B_b$: the class
$\mathcal H=\{f\in B_b(\C): P_tf\in C_b(\C)\}$ is a vector space, contains
$C_b(\C)$, and is closed under bounded pointwise monotone limits (dominated
convergence gives $P_tf_k\to P_tf$ pointwise, and the uniform Lipschitz bound
\eqref{eq:limitbound} — which depends only on $\|f\|_\infty$ — is preserved,
keeping the limit in $C_b$). By the functional monotone class theorem,
$\mathcal H=B_b(\C)$, since $C_b(\C)$ generates the Borel $\sigma$-algebra.
Therefore $P_t$ is strong Feller.
\end{proof}

\begin{remark}
The monotone-class extension genuinely uses that the Lipschitz constant in
\eqref{eq:limitbound} depends on $f$ only through $\|f\|_\infty$ (not through any
modulus of continuity of $f$). This is the typical strength of gradient
estimates of BEL type and is exactly what Lemma~\ref{lem:unifgrad} provides.
\end{remark}

\section{Conclusion: uniqueness of the invariant measure}

\begin{theorem}\label{thm:main}
The Markov semigroup $P_t$ on $\C$ has at most one invariant probability
measure.
\end{theorem}

\begin{proof}
Fix any $t_0>0$. By Proposition~\ref{prop:sf}, $P_{t_0}$ is strong Feller on the
Polish space $\C$. By Proposition~\ref{prop:irr}, $P_{t_0}$ is irreducible:
$P_{t_0}\mathbf 1_U(x)>0$ for every nonempty open $U$ and every $x\in\C$. By the
Doob--Khasminskii theorem (Theorem~\ref{thm:doob}), strong Feller plus
irreducibility imply that the transition kernels $\{P_{t_0}(x,\cdot):x\in\C\}$
are mutually equivalent, and hence $(P_t)$ admits at most one invariant
probability measure.

Concretely: if $\mu,\nu$ were both invariant, then for the $P_{t_0}$-stationary
measures we would have, for every Borel $A$,
$\mu(A)=\int P_{t_0}(x,A)\,\mu(dx)$ and likewise for $\nu$. Suppose
$\mu\ne\nu$; then there is a Borel set $A$ with, say, $\mu(A)=1$, $\nu(A)<1$
(two distinct invariant measures of a Markov semigroup with the equivalence
property are mutually singular if distinct — Doob's theorem — yet equivalence
forces them to be mutually absolutely continuous, a contradiction). Hence
$\mu=\nu$.
\end{proof}

\section{Summary of the argument and where the singular drift was used}

\begin{enumerate}
\item \textbf{Interpretation.} $\delta_0(u)$ is the occupation-density (local
time) drift; its regularisations $G_{1/n}$ are nonnegative probability densities
of total mass $1$, uniformly in $n$. The additive space--time white noise on the
bounded interval $[0,1]$ with Dirichlet boundary conditions is the source of all
the ergodic regularity.

\item \textbf{Irreducibility} (Prop.~\ref{prop:irr}) follows from the full
support of the Dirichlet stochastic convolution $Z(t)$ on $\C$, together with the
\emph{nonnegativity and finite mass} of the limiting drift shift $\Phi(t)\ge0$;
a fixed-sign bounded shift cannot kill the full Gaussian support.

\item \textbf{Strong Feller} (Prop.~\ref{prop:sf}) follows from a
Bismut--Elworthy--Li gradient estimate \eqref{eq:gradbound} for the regularised
semigroups that is \emph{uniform in $n$} because it depends on the drift only
through $\|G_{1/n}\|_{L^1}=1$, not through $\|G_{1/n}\|_\infty\to\infty$. The
uniform Lipschitz bound passes to the limiting $P_t$ via the assumed weak
convergence (H2) and a monotone-class argument.

\item \textbf{Conclusion.} Strong Feller $+$ irreducibility $\Rightarrow$ at
most one invariant measure, by Doob--Khasminskii. \emph{Hence the statement is
proved: $P_t$ has at most one invariant probability measure.}
\end{enumerate}

\paragraph{Honest assessment.} The proof is complete and rigorous \emph{except}
for the $n$-uniform gradient estimate of Lemma~\ref{lem:unifgrad}
(Remark~\ref{rmk:gap}), whose precise hypotheses I state and whose mechanism
(total-mass control of the singular drift, the SPDE analogue of the bounded
local-time Girsanov exponent for skew Brownian motion) I justify, but whose full
proof requires the stochastic-sewing/Malliavin estimates of [ABLM24] adapted to
the gradient of the semigroup. Modulo that single, clearly identified analytic
input, the answer to the question is: \textbf{the statement is TRUE; $P_t$ has at
most one invariant probability measure.}

\end{document}
